
\scanpage{61}
\mspage{17}
Récapitulation des opérations élémentaires

\inserted{$(0^\circ)$ Somme disjointe}

$1^\circ$) \uline{gommage} \inserted{\uline{total} des repères en $a$, i.e.} des repères se trouvant sur un cercle de recollement \inserted{$C_a$}, ou \deleted{\unsure{sur un cercle}} de découpage
$\Longleftrightarrow$ suppression de la \deleted{partie} \inserted{$R_a$}
$(a\in A)$ $\Longleftrightarrow$ modification de $r_a$ en le
\deleted{\unsure{de $R$ se trouvant dans une partie}}
remplaçant par $0$ en $a$.
\margin{\uline{fortement irréversible}\quad $r_a\ne0$}

$2^\circ$) \inserted{\uline{gommage partiel uniforme}: \uline{d'exposant $d_a\mid r_a$}}
Restriction des \deleted{scalaires} \inserted{groupes de $R_a$} de $\mathbb Z_a/r_a\mathbb Z$ à un
$\mathbb Z_a/r'_a\mathbb Z_a$, où $r'_a\mid r_a$
(pour $a$ fixé, il y a $r_a/r'_a=d_a$ choix possibles, si on suppose $d_a>1$ i.e. $r'_a\ne r_a$).
\margin{\uline{renforcement de structure faible}\quad $r_a\unsure{\ge|>}1$
\[
r'_a\unsure{\mid|\ne}r_a\text{ et }r'_a\ne r_a
\qquad\text{i.e. }d_a=\frac{r_a}{r'_a}\ne1.
\]}

$3^\circ$) \deleted{Extension des}
\inserted{\uline{Surmarquage uniforme} \uline{d'exposant} $d_a>1$}
--- i.e. extension des groupes de $R_a$ de
$\mathbb Z_a/r_a\mathbb Z_a$ à $\mathbb Z_a/d_ar_a\mathbb Z_a$,
\margin{\uline{affaiblissement de structure faible}\quad $r_a\ne0$}

\textcircled{$4^\circ$} gommage du cercle de découpage $C_a$ $(a\in A_0)$

\textcircled{$5^\circ$} bouchage du trou $C_a$ $(a\in A_l)$
\margin{\uline{fortement irréversible}\quad $\left\{\begin{array}{l}4^\circ\\5^\circ\end{array}\right.$}

$6^\circ$) Découpage le long de $C_a$ $(a\in A_0)$ \({}^*\)
\margin{fortement irréversible si $r_a=0$, \deleted{faible} réversible = anodin si $r_a\unsure{\ge|>}1$.}

$7^\circ$) Recollement le long de $C_a,C_{a'}$
(avec $a,a'\in A_l$, $a\ne a'$, \inserted{$r_a=r_{a'}$}), via un élément de
$R_a\mathbin{\wedge}R_{a'}$.
\margin{$r_a=r_{a'}\unsure{\ge|>}1$\quad réversible = anodin.}

\margin{\({}^*\) Il faudrait \inserted{peut-être} séparer $6^\circ)$ en deux, suivant que $r_a=0$ ou $r_a\ne0$ (ou exclure le cas $r_a=0$?).}
\mspage{60}
